\section{Introduction}
\label{sec:introduction}

Finite-set dualities turn local observables of an interacting particle
system into observables of a process on finite subsets of the lattice.
For two-state spin systems, the ordinary monomials
\[
  \chi_A(\eta)=\prod_{i\in A}\eta(i),
  \qquad A\Subset\Lambda,
\]
give such a duality without requiring attractiveness or a special
algebraic structure in the flip rates.  The price is that the dual
process is signed and carries a Feynman--Kac weight.  Its individual
transitions therefore do not yield the positive pathwise
representations that make classical set-valued dualities useful.

This paper recovers positivity after a different conditioning.  We
reveal only those nonempty-target dual interactions whose source is
active.  These successful interactions cut each site-line into
spacetime intervals, which we call patches.  Once their locations and
targets are fixed, the interaction histories inside distinct patches
are independent after imposing an elementary local consistency event.
The signed Feynman--Kac weight can then be averaged one patch at a
time.  This produces an exact random product whose factors are
one-site quantities.

The construction applies to general finite-range spin systems.  No
orientation of the lattice is assumed.  A uniform \(1\)-to-\(0\)
component appears only in the convergence theorem of
\cref{sec:common-limit}; it is not part of the representation or
positivity theory.

\subsection{Main results}

\paragraph{Patch representation.}
The monomial duality of \cref{sec:monomial-duality} is refined in
\cref{sec:patches,sec:representation}.  Conditional on the
finite-horizon successful-interaction sigma algebra, all patch
histories factorize.  If \(\cB_t\) and \(\cE_t\) denote the bulk and end
patches at time \(t\), then
\begin{equation}
  P_t\chi_A(\xi)
  =
  \E_A\left[
    \prod_{P\in\cB_t}C(P)
    \prod_{P\in\cE_t}C(\xi(i(P)),P)
  \right].
  \label{eq:intro-representation}
\end{equation}
The factor \(C(P)\) depends only on the label, duration, source site,
and initial target of \(P\).  End factors are affine in the terminal
spin.  The factorization theorem is finite-horizon; no conditioning on
the all-time skeleton is needed.

\paragraph{Patch positivity.}
The explicit one-patch calculation leads to a local coefficient
criterion, stated in \cref{thm:patch-positivity-criterion}, that is
necessary and sufficient for all bulk factors in
\eqref{eq:intro-representation} to be nonnegative.  The same
calculation determines the smallest terminal one-density profile
\(\pstar\) for which all end factors are nonnegative:
\begin{equation}
  p_i^\star
  =
  \max\left\{
    0,
    \sup_{\substack{\vn\ne S\subseteq N(i)\\
                    c_i^0(S)+c_i^1(S)<0}}
    \frac{c_i^0(S)}{c_i^0(S)+c_i^1(S)}
  \right\}.
  \label{eq:intro-critical-density}
\end{equation}
This is a representation threshold, not a thermodynamic critical
density.

\paragraph{Centered-moment order.}
The profile \(\pstar\) defines centered monomials
\[
  \chi_A^\star(\eta)
  =
  \prod_{i\in A}\bigl(\eta(i)-p_i^\star\bigr)
\]
and the high-density cone
\[
  \Mstar
  =
  \set{\mu:\mu(\chi_A^\star)\ge0
  \text{ for every }A\Subset\Lambda}.
\]
In \cref{sec:centered-order} we show that a patch-positive semigroup
preserves the order induced by all centered moments.  In particular,
\(\Mstar\) is invariant under the dynamics.  This conclusion is an
order of joint occupation moments and need not imply stochastic
domination.

\paragraph{Comparison under removal of pure deaths.}
Suppose
\[
  \cL=\cL'+\sum_{i\in\Lambda}
  a_i\bigl(f(\eta^{i,0})-f(\eta)\bigr),
  \qquad a_i\ge0.
\]
If \(\cL\) is patch-positive, then so is \(\cL'\), with the same
critical profile, and
\[
  \nu(P_t\chi_A)\le \nu(P_t'\chi_A),
  \qquad \nu\in\Mstar.
\]
The proof in \cref{sec:pure-death-comparison} compares the positive
random products skeleton by skeleton.

\paragraph{A common invariant limit.}
Assume in addition that the lattice has polynomial volume growth of
exponent \(D\) and
\[
  c_i^1(\eta)\ge\varepsilon>0
  \qquad(i\in\Lambda,\ \eta\in\{0,1\}^{\Lambda}).
\]
Then there is an invariant probability measure \(\pi\) such that, for
every local \(f\),
\begin{equation}
  \sup_{\nu\in\Mminus}
  \abs{\nu(P_tf)-\pi(f)}
  \le K_f(1+t)^D e^{-\varepsilon t/2}.
  \label{eq:intro-common-limit}
\end{equation}
Here
\[
  \Mminus
  =
  \set{\mu:
    \mu(\chi_A^\star)\ge-\chi_A^\star(\mathbf1)
    \text{ for every }A\Subset\Lambda}
\]
is larger than \(\Mstar\).  The limiting monomial moments have an
explicit product formula over the full patch family.  If
\(\pstar\le\frac12\mathbf1\), then every probability measure belongs
to \(\Mminus\), and \eqref{eq:intro-common-limit} becomes uniform
exponential ergodicity over all initial configurations.

\subsection{Examples and scope}

The FA-1f model and the biased annihilating branching process both
satisfy patch positivity, and in both cases \(p^\star\) equals the
equilibrium one-density.  Soft FA-1f illustrates the common-limit
theorem.  For the unperturbed biased annihilating branching process,
whose rates vanish at the all-one configuration, centered-moment order
instead extends known homogeneous convergence to a class of
inhomogeneous and non-product high-density initial laws.

The signed monomial duality itself belongs to the general theory of
Markov-process duality and additive or cancellative graphical
constructions; see, for example,
\citet{liggett1985,jansen-kurt2014,sudbury-lloyd1995}.  The contribution
here is the successful-interaction conditioning, its one-site patch
factorization, and the consequences obtained from positivity of the
averaged patch factors.

\subsection{Organization}

\Cref{sec:monomial-duality} derives the signed monomial dual.
\Cref{sec:patches} constructs patches and proves conditional
factorization.  \Cref{sec:representation} gives the random-product
representation.  \Cref{sec:patch-positivity} identifies the local
positivity criterion, the critical profile, and the centered cones.
\Cref{sec:centered-order,sec:pure-death-comparison} prove the order and
comparison results.  \Cref{sec:common-limit} proves the common-limit
theorem, and \cref{sec:examples} treats FA-1f and BABP.  The local patch
calculation and the confinement lemmas are collected in the
appendices.
